%% LyX 2.6.0~devel created this file.  For more info, see https://www.lyx.org/.
%% Do not edit unless you really know what you are doing.
\documentclass[12pt,english,hebrew]{article}
\usepackage{amsmath}
\usepackage{fontspec}
\usepackage{unicode-math}
\usepackage[a4paper]{geometry}
\geometry{verbose,tmargin=2cm,bmargin=2cm,lmargin=1cm,rmargin=1cm}
\usepackage{esint}
\usepackage{microtype}

\makeatletter

%%%%%%%%%%%%%%%%%%%%%%%%%%%%%% LyX specific LaTeX commands.
\providecommand\textquotedblplain{%
  \bgroup\addfontfeatures{Mapping=}\char34\egroup}

%%%%%%%%%%%%%%%%%%%%%%%%%%%%%% User specified LaTeX commands.
% -------------------------------------------------
% Encoding and Fonts (fontspec is usually handled by LyX with XeTeX)
% \usepackage{fontspec} % LyX will add this if XeTeX/LuaTeX is used.
% \usepackage{polyglossia} % LyX handles this with language settings.
% \setmainlanguage{hebrew} % Set this in LyX: Document > Settings > Language
% \setotherlanguage{english} % Set this in LyX: Document > Settings > Language

\newfontfamily\hebrewfont{David CLM} % Keep for custom commands/styles
\newfontfamily\hebrewfontsf{Nachlieli CLM} % Keep for custom commands/styles
\newfontfamily\englishfont{Latin Modern Roman} % Keep for custom commands/styles

% -------------------------------------------------
% Math and Symbols
\usepackage{amsmath, mathtools, amsthm} % amsmath & amsthm often by LyX, mathtools is a good addition.
% Add this if not automatically handled by LyX's math font settings.
% \setmathfont{Latin Modern Math} % Set this in LyX: Document > Settings > Math Options > Use unicode-math
                                % If you set it here, it overrides LyX's GUI setting.
\setmathfont{XITS Math}
%\setmathfont{Libertinus Math}
\usepackage{physics}          % For \bra, \ket, etc.
\usepackage{autobreak}        % Break long equations automatically
\usepackage{relsize}          % For \mathlarger
\usepackage{microtype}        % Better text rendering (often default with fontspec)

% -------------------------------------------------
% Document Layout
% \usepackage{geometry} % Configure in LyX: Document > Settings > Page Layout
% \geometry{a4paper, margin=1in} % Configure in LyX
% \usepackage{setspace} % Configure in LyX: Document > Settings > Text Layout
% \onehalfspacing % Configure in LyX
\usepackage{float}            % Improved float control

% -------------------------------------------------
% Titles
\usepackage{titlesec}
\titleformat{\section}{\Large\bfseries\hebrewfont}{\thesection}{10pt}{}
\titleformat{\subsection}{\large\bfseries\hebrewfont}{\thesubsection}{10pt}{}
\titleformat{\subsubsection}{\normalsize\bfseries\hebrewfont}{\thesubsubsection}{10pt}{}

% -------------------------------------------------
% Headers
\usepackage{fancyhdr}
\pagestyle{fancy}
\fancyhead{}
\fancyhead[R]{\leftmark} % Or \rightmark, depending on desired output with polyglossia/Hebrew
\renewcommand{\headrulewidth}{0.5pt}
% \setlength{\footskip}{12pt} % If needed, keep. Default might be okay. Avoids fancyhdr warning.

% -------------------------------------------------
% Lists
\usepackage{enumitem}
\setlist[itemize]{label=\textbullet, leftmargin=1.5em, nosep}
\setlist[enumerate]{leftmargin=1.5em, labelsep=0.5em, nosep}

% -------------------------------------------------
% Footnote Rule
\renewcommand{\footnoterule}{%
  \kern 3pt
  \hbox to \textwidth{\hfill\vrule height 0.5pt width 0.4\textwidth}
  \kern 4pt
}

% -------------------------------------------------
% Math Display Spacing
\AtBeginDocument{
  \setlength\abovedisplayskip{6pt}
  \setlength\belowdisplayskip{6pt}
  \setlength\abovedisplayshortskip{6pt}
  \setlength\belowdisplayshortskip{6pt}
  \relpenalty=9999
  \binoppenalty=9999
}

% -------------------------------------------------
% RTL Math Tag Fix
\makeatletter
\def\maketag@@@#1{\hbox{\m@th\normalfont\LRE{#1}}}
\def\tagform@#1{\maketag@@@{(\ignorespaces#1\unskip)}}
\makeatother

% -------------------------------------------------
% QED Symbol
\renewcommand{\qedsymbol}{$\blacksquare$} % Requires amsthm

% -------------------------------------------------
% Hebrew inside math mode
% \usepackage{amsmath} % Already loaded above, \text command source
\newcommand{\hebtext}[1]{\text{\hebrewfont #1}}

\makeatother

\usepackage{polyglossia}
\setdefaultlanguage{hebrew}
\setotherlanguage{english}
\begin{document}
$\date{14/04/2026}$
\begin{english}[variant=american]%

\section*{Form factor}

\end{english}%
ניקח אלומת אלקטרונים ונקרין אותם על הגרעין.

לאלקטרון אורך גל קטן.

נרצה להשתמש באלקטרונים בעלי אורך גל קטנים משל אורך הכל של הגרעין.
(בערך $\sim100Mev$).

האלקטרון אדיש לכוח הגרעיני החזק.

ניתן לקבל מידע על התפלגות המסה של הגרעין ע\textquotedbl י אופי הפיזור
של האלקטרון מהגרעין.

\[
f\left(k-k'\right)=\intop\psi^{\ast}V\psi d^{3}r
\]

כאשר $\psi=e^{-ik\cdot r},\qquad\psi'=e^{-ik'\cdot r}$
\[
f=\intop e^{i\overset{q}{\overbrace{\left(k-k'\right)}}\cdot r}\cdot Vd^{3}r
\]
\begin{equation}
F\left(q\right)=\frac{1}{Ze}\intop d^{3}r\rho\left(r\right)e^{iq\cdot r}\label{eq:Scattering Amplitude}
\end{equation}

$F\left(q\right)$ היא אמפליטודת הפיזור (\textenglish[variant=american]{(Form
Factor}

צפיפות הגרעינים של אטומים שונים לא שונה במיוחד.

בערך $\sim0.15\frac{n}{fm^{3}}$ זוהי גם הצפיפות של כוכב נויטרונים.

רדיוס- $R=R_{0}A^{1/3}$

כאשר $R_{0}=1.25\ fm$ .

\paragraph{סימטריית איזו-ספין}

\[
^{A}_{Z}X_{N}\leftrightarrow{}^{A}_{N-1}X_{Z+1}
\]

\[
E=\frac{3}{5}\frac{Z^{2}e^{2}}{4\pi\varepsilon_{0}R}
\]

\begin{english}[variant=american]%

\section*{Mass and binding energy}

\end{english}%
\[
M_{nuc}\left(A,Z\right)=M_{atom}\left(A,Z\right)-Zm_{e}
\]
 כאשר מזניחים את אנרגיית הקשר האלקטרונית.

\subsubsection*{אנרגיית הקשר הגרעינית ניתנת לחישוב ע\textquotedbl י 
\[
B\left(A,Z\right)=\left[Zm_{p}+Nm_{n}-M_{nuc}\left(A,Z\right)\right]c^{2}
\]
\[
B\left(A,Z\right)=\left[Zm_{p}+Zm_{e}+Nm_{n}-M_{nuc}\left(A,Z\right)-M_{nuc}\left(A,Z\right)-Z\cdot m_{e}\right]c^{2}
\]
\[
B\left(A,Z\right)=\left[Z\cdot m_{H}+Nm_{n}-M_{atom}\left(A,Z\right)\right]c^{2}
\]
}

\paragraph*{יחידות: }

יחידות מסה אטומית $a.m.u=u=931.5MeV/c^{2}$ 

\paragraph*{פיחות המסה \textenglish[variant=american]{Mass Defect }}

\[
\Delta\equiv\left(Au-m\right)
\]

לדוגמה 
\[
m_{17}=7.016...u\Rightarrow\Delta=0.016...u
\]

\begin{english}[variant=american]%

\subsubsection*{Q-Value\texthebrew{ }}

\end{english}%
\[
x\rightarrow y
\]
\[
Q=\left(m_{x}-m_{y}\right)c^{2}
\]
\[
^{8}Be\rightarrow He+He
\]
\[
Q=\left(m_{^{8}Be}-2m_{He}\right)c^{2}
\]
 $Q>0$ - תהליך אקזו-טרמי שיכול להתקיים ספונטינית בטבע.

$Q<0$ - תהליך אנדו-טרמי שצריך להשקיע אנרגיה כדי לקיימו.
\begin{english}[variant=american]%

\subsection*{Neutron-Proton seperation energy}

\end{english}%
\begin{align*}
S_{n}\left(A,Z\right) & =\left[\underset{after}{\underbrace{M\left(A-1,Z\right)}}+m_{n}-\underset{before}{\underbrace{M\left(A,Z\right)}}\right]c^{2}\\
S_{p}\left(A,Z\right) & =\left[M\left(A-1,Z-1\right)+m_{H}-M\left(A,Z\right)\right]c^{2}
\end{align*}


\subsubsection*{ספקטרוסקופיית מסות}

\[
\frac{m}{q}=r\frac{B^{2}}{E}
\]

\begin{english}[variant=american]%

\subsubsection*{Multi reflection time of flight mass spectrometer }

\end{english}%
\[
qV=\frac{1}{2}mv^{2}
\]
\[
v=\sqrt{\frac{2qV}{m}}
\]
 
\begin{english}[variant=american]%

\subsection*{Penning Trap }

\end{english}%
שדה מגנטי ושדה קוואדרופולי חשמלי . החלקיק נלכד ועושה סיבובים.
\end{document}
